\begin{alist}
\item Από τη γραφική παράσταση της $f$ συμπεραίνουμε ότι αυτή είναι γνησίως αύξουσα στο διάστημα $[-1, 0]$ και γνησίως φθίνουσα στο $[0, +\infty)$. Επιπλέον η $f$ παρουσιάζει ολικό μέγιστο για $x = 0$. Η μέγιστη τιμή της είναι ίση με $f(0) = 1$.

\item Οι αριθμοί $-\dfrac{3}{5}, -\dfrac{5}{9}$ περιέχονται στο διάστημα $[-1, 0]$ όπου η συνάρτηση είναι γνησίως αύξουσα και
\[-\frac{3}{5} + \frac{5}{9} = \frac{-27+25}{45} = -\frac{2}{45} < 0\]
οπότε $-\dfrac{3}{5} < -\dfrac{5}{9}$ και λόγω της μονοτονίας της $f$ συμπεραίνουμε ότι $f\left(-\dfrac{3}{5}\right) < f\left(-\dfrac{5}{9}\right)$.

\item Είναι:
\[\sigma\upsilon\nu 120^\circ = \sigma\upsilon\nu(180^\circ - 60^\circ) = -\sigma\upsilon\nu 60^\circ = -\frac{1}{2} \in [-1, 0]\]
και 
\[\eta\mu 120^\circ = \eta\mu(180^\circ - 60^\circ) = \eta\mu 60^\circ = \frac{\sqrt{3}}{2} \in (0, +\infty).\]
Έτσι, με τη βοήθεια του τύπου της $f$ έχουμε:
\[f(\sigma\upsilon\nu 120^\circ) = f\left(-\frac{1}{2}\right) = \sqrt{1 - \frac{1}{4}} = \frac{\sqrt{3}}{2}\]
και 
\[f(\eta\mu 120^\circ) = f\left(\frac{\sqrt{3}}{2}\right) = 1 - \frac{\sqrt{3}}{2} = \frac{2 - \sqrt{3}}{2}.\]

\item Η γραφική παράσταση της $g$ προκύπτει με μεταφορά της $\text{C}_f$ δυο μονάδες προς τα δεξιά και φαίνεται στο επόμενο σχήμα.
\end{alist}
\begin{center}
\begin{tikzpicture}
\begin{axis}[
    axis lines=middle,
    xmin=-1.5, xmax=4.5,
    ymin=-1.5, ymax=1.5,
    xtick={-1, -0.5, 0, 0.5, 1, 1.5, 2, 2.5, 3, 3.5, 4},
    ytick={-1, -0.5, 0.5, 1},
    grid=none,
    thick,
    width=14cm, height=8cm,
    samples=100
]
\addplot[domain=-1:0, black, thick] {sqrt(1 - x^2)};
\addplot[domain=0:1.5, black, thick] {1 - x};
\draw[black, thick] (axis cs:-1,0) -- (axis cs:0,1);
\node at (axis cs:-0.7, 0.8) {$\text{C}_f$};

\addplot[domain=1:2, black, thick, dotted] {sqrt(1 - (x-2)^2)};
\addplot[domain=2:3.5, black, thick, dotted] {1 - (x-2)};
\draw[black, thick, dotted] (axis cs:1,0) -- (axis cs:2,1);
\node at (axis cs:1.7, 1.1) {$\text{C}_g$};
\end{axis}
\end{tikzpicture}
\end{center}
